% ============ 第3章 §3.1 交换Banach代数(原书 pp.203-226) ============
\section{交换 Banach 代数}\label{sec:3.1}
\index{Banach 代数}
\index{Banach 代数!交换----}
\index{Banach 代数!添加单位元的交换----}
\index{交换 Banach 代数}
\index{谱}
\index{预解式}
\index{谱半径}
\index{理想}
\index{理想!极大----}
\index{理想!正则----}
\index{理想!闭----}
\index{正则!(理想)}
\index{Gelfand 变换}
\index{Gelfand 表示}
\index{极大理想空间}
\index{范数拓扑}
\index{*-弱拓扑}
\index{hk 拓扑}
\index{谱综合}
\index{谱综合集}
\index{谱综合!谱综合可能}
\index{谱综合!谱综合不可能}
\index{半单}
\index{闭图像定理}
\index{特征}
\index{Ditkin 条件}
\index{作用函数!(半单交换 Banach 代的)}


本节的目的在于叙述交换 Banach 代数的一般理论.Fourier 分析中常出现的交换 Banach 代数,其极大理想空间的结构往往复杂,但这里的例子处理的是简单的情形.下一节将处理 Fourier 分析中出现的主要交换 Banach 代数.另外,函数分析中众所周知的基础知识视为已知.

\subsection{交换 Banach 代数的定义}\label{sec:3.1.1}

\begin{definition}\label{def:3.1.1}
集合 $A$ 具有如下性质.

(1) $A$ 是复数域 $\mathbf{C}$ 上的环 $(+, \cdot)$.即,

(i) $A$ 是 $\mathbf{C}$ 上的向量空间.

(ii) $x(yz) = (xy)z$,$x(y + z) = xy + xz$,$\alpha(xy) = (\alpha x)y = x(\alpha y)$($x, y, z \in A$,$\alpha \in \mathbf{C}$).

(2) 对 $x \in A$,赋予非负值 $\|x\|$ 使其满足下列条件(把满足下述 (i), (ii), (iii) 的 $\|\cdot\|$ 称为范数).

(i) $\|x\| \geq 0$,且 $\|x\| = 0$ 的情形仅限于 $x = 0$.

(ii) $\|\alpha x\| = |\alpha|\,\|x\|$.

(iii) $\|x + y\| \leq \|x\| + \|y\|$.

(vi) $\|xy\| \leq \|x\| \cdot \|y\|$.

(3) $A$ 关于由范数导出的距离 $d(x, y) = \|x - y\|$ 完备.

此时,称 $A$ 为 Banach 代数.特别地,当对 $x, y \in A$ 有 $xy = yx$ 时,称 $A$ 为交换 Banach 代数.又,当存在满足 $ex = xe = x$ 的 $e \in A$ 时,称 $e$ 为 $A$ 的单位元.以后,$A$ 中的拓扑均按由范数 $\|\cdot\|$ 导出的距离 $d(x, y)$ 来考虑.
\end{definition}

在交换 Banach 代数中,若存在单位元 $e$,则设 $\|e\| = 1$.若不存在单位元,令
\begin{equation*}
A_{e} = \{(x, \lambda):\ x \in A,\ \lambda \in \mathbf{C}\}.
\end{equation*}
此时,定义
和 $(x, \lambda) + (y, \mu) = (x + y, \lambda + \mu)$,
数乘 $\alpha(x, \lambda) = (\alpha x, \alpha\lambda)$,
积 $(x, \lambda) \cdot (y, \mu) = (xy + \mu x + \lambda y, \lambda\mu)$,
\begin{align*}
\text{范数} & \|(x, \lambda)\| = \|x\| + |\lambda|, \\
& (x, y \in A,\ \lambda, \mu, \alpha \in \mathbf{C})
\end{align*}
则 $A_{e}$ 成为以 $e = (0, 1)$ 为单位元的交换 Banach 代数.此时,在 $A_{e}$ 中,$\{(x, 0):\ x \in A\}$ 作为交换 Banach 代数可与 $A$ 等同.称这个 $A_{e}$ 为对 $A$ 添加单位元而得到的交换 Banach 代数.

例．(1) 若 $\mathbf{T} = \{|z| = 1:\ z \in \mathbf{C}\}$,则 $\mathbf{T}$ 是紧集,且关于乘法构成交换群.又,若以 $C(\mathbf{T})$ 表示 $\mathbf{T}$ 上复值连续函数的全体,则按通常的函数运算与范数 $\|f\|_{\infty} = \sup\{|f(z)|:\ z \in \mathbf{T}\}$,$C(\mathbf{T})$ 成为交换 Banach 代数.

(2) 令 $\ell^{1} = \{a = \{a_{n}\}:\ \sum_{n \in \mathbf{Z}} |a_{n}| < \infty,\ a_{n} \in \mathbf{C}\ (n \in \mathbf{Z})\}$.$\mathbf{Z}$ 表示整数集.又,对 $a = \{a_{n}\}$,$b = \{b_{n}\} \in \ell^{1}$,$\alpha, \beta \in \mathbf{C}$,有 $\alpha a + \beta b \in \ell^{1}$.进一步,令 $\|a\|_{\ell^{1}} = \sum_{n \in \mathbf{Z}} |a_{n}| < \infty$,$a * b = \{\sum_{k \in \mathbf{Z}} a_{n-k} b_{k}\}$,则 $\ell^{1}$ 关于范数 $\|\cdot\|_{\ell^{1}}$ 与运算 $*$ 成为交换 Banach 代数.

此后,对紧 Hausdorff 空间 $X$,用 $C(X)$ 表示以 $\|f\|_{\infty} = \sup_{x \in X} |f(x)|$ 为范数的 Banach 空间.

\subsection{逆元,预解式与谱}\label{sec:3.1.2}

\begin{definition}\label{def:3.1.2}
设 $A$ 为具有单位元 $e$ 的交换 Banach 代数,令
\begin{equation*}
A^{-1} = \{x \in A:\ \text{存在 } x \text{ 的逆元 } x^{-1} \text{ 使得 } xx^{-1} = e\}.
\end{equation*}
\end{definition}

\begin{theorem}\label{thm:3.1.3}
设 $A$ 为具有单位元 $e$ 的交换 Banach 代数.

(1) $A^{-1} \supset \{x:\ \|e - x\| < 1\}$,且 $A^{-1}$ 是开集.

(2) 由 $x$ 到 $x^{-1}$ 的对应是从 $A^{-1}$ 到 $A^{-1}$ 的连续映射.
\end{theorem}

\begin{proof}
(1) 当 $\|x\| < 1$ 时,令 $x^{0} = e$,则 $\sum_{n=0}^{\infty} x^{n}$ 在 $A$ 中绝对收敛.因此 $\sum_{n=0}^{\infty} x^{n}$ 收敛.此时
\begin{equation*}
(e - x)\sum_{n=0}^{\infty} x^{n} = \sum_{n=0}^{\infty} x^{n} - \sum_{n=1}^{\infty} x^{n} = e
\end{equation*}
故 $\{e - x:\ \|x\| < 1\} \subset A^{-1}$.于是 $\{x:\ \|e - x\| < 1\} \subset A^{-1}$.

进一步,当 $x \in A^{-1}$ 时,若 $\|y\| < \|x^{-1}\|^{-1}$,则
\begin{equation*}
\|x^{-1} y\| \leq \|x^{-1}\| \cdot \|y\| < 1
\end{equation*}
故由前半可知 $e - (-x^{-1}y) \in A^{-1}$.从而若 $\|y\| < \|x^{-1}\|^{-1}$,则 $x + y \in A^{-1}$,可见 $A^{-1}$ 是开集.

(2) 利用 (1),
\begin{align*}
(x + y)^{-1} - x^{-1} & = x^{-1}\bigl((e + x^{-1}y)^{-1} - e\bigr) \\
& = x^{-1}\Bigl(\sum_{n=0}^{\infty} (-x^{-1}y)^{n} - e\Bigr) \\
& = x^{-2}y\sum_{n=0}^{\infty} (-1)^{n+1}(x^{-1}y)^{n},
\end{align*}
故
\begin{align*}
\|(x + y)^{-1} - x^{-1}\| & \leq \|x^{-2}y\|\sum_{n=0}^{\infty}\|x^{-1}y\|^{n} \\
& \leq \|x^{-2}\| \cdot \|y\|\sum_{n=0}^{\infty}\bigl(\|x^{-1}\| \cdot \|y\|\bigr)^{n} \\
& = \frac{\|x^{-2}\| \cdot \|y\|}{1 - \|x^{-1}\| \cdot \|y\|} \to 0 \quad (\|y\| \to 0)
\end{align*}
由此可知,当 $\|y\| \to 0$ 时,$\|(x + y)^{-1} - x^{-1}\| \to 0$.
\end{proof}

\begin{definition}\label{def:3.1.4}

(1) 设 $A$ 为具有单位元 $e$ 的交换 Banach 代数.

(i) 对 $x \in A$,
\begin{equation*}
\Sp_{A}(x) = \{\lambda \in \mathbf{C}:\ x - \lambda e \notin A^{-1}\}
\end{equation*}
称为 $x$ 的(在 $A$ 中的)谱.$\Sp_{A}(x)$ 也记作 $\Sp(x, A)$.

(ii) 对满足 $\lambda \notin \Sp_{A}(x)$ 的 $\lambda \in \mathbf{C}$,
\begin{equation*}
R(x, \lambda) = (x - \lambda e)^{-1}
\end{equation*}
称为 $x$ 的预解式.

(2) 当 $A$ 为没有单位元的交换 Banach 代数时,令 $\Sp_{A}(x) = \Sp_{A_{e}}(x)$.此时容易看出,对 $\lambda \neq 0$,$\lambda \notin \Sp_{A}(x)$ 的充要条件是存在 $y \in A$ 使得 $xy - \lambda^{-1}x - \lambda y = 0$.
\end{definition}

\begin{definition}\label{def:3.1.5}
设 $\mathbf{C} \supset D$ 为开集,对交换 Banach 代数 $A$,若对任意 $\lambda_{0} \in D$,
\begin{equation*}
\lim_{\lambda \to \lambda_{0}} \frac{F(\lambda) - F(\lambda_{0})}{\lambda - \lambda_{0}}
\end{equation*}
在 $A$ 中存在,则称 $F(\lambda):\ D \to A$ 在 $D$ 中正则.
\end{definition}

\begin{theorem}\label{thm:3.1.6}
设 $A$ 为交换 Banach 代数.此时,对 $x \in A$,下述 (i)～(iii) 成立.

(i) $\Sp_{A}(x)$ 是闭集.

(ii) $\Sp_{A}(x) \subset \{\lambda:\ |\lambda| \leq \|x\|\}$.

(iii) $\Sp_{A}(x) \neq \emptyset$.

特别地,由 (i), (ii) 可知 $\Sp_{A}(x)$ 是复平面上的紧集.
\end{theorem}

\begin{proof}
首先设 $A$ 具有单位元 $e$.(i) 令 $F(\lambda) = x - \lambda e$,则 $F(\lambda):\ \mathbf{C} \to A$ 是连续映射.事实上,由
\begin{equation*}
\|F(\lambda) - F(\lambda_{0})\| = \|(x - \lambda e) - (x - \lambda_{0}e)\| = |\lambda - \lambda_{0}|\,\|e\|
\end{equation*}
可知.这里,由定理~\ref{thm:3.1.3},$A^{-1}$ 是开集,故 $(A^{-1})^{c}$ 是闭集.于是因 $F$ 连续,$F^{-1}((A^{-1})^{c})$ 是复平面的闭子集.从而 $\Sp_{A}(x) = F^{-1}((A^{-1})^{c})$ 是闭集.

(ii) 若 $|\lambda| > \|x\|$,则由定理~\ref{thm:3.1.3},$x - \lambda e = -\lambda(e - \lambda^{-1}x) \in A^{-1}$,故 $\lambda \notin \Sp_{A}(x)$.因此 $\Sp_{A}(x) \subset \{\lambda:\ |\lambda| \leq \|x\|\}$.从而 $\Sp_{A}(x)$ 是 $\mathbf{C}$ 的紧子集.

(iii) 证明 $\Sp_{A}(x) \neq \emptyset$.当 $\lambda_{0} \notin \Sp_{A}(x)$ 时,$G(\lambda) = R(x, \lambda)$ 在 $\lambda = \lambda_{0}$ 处正则,且当 $|\lambda| > \|x\|$ 时,
\begin{equation*}
\lim_{\lambda \to \lambda_{0}} \frac{R(x, \lambda) - R(x, \lambda_{0})}{\lambda - \lambda_{0}} = \bigl(R(x, \lambda_{0})\bigr)^{2},
\end{equation*}
\begin{equation*}
R(x, \lambda) = -\frac{1}{\lambda}\sum_{n=0}^{\infty}\left(\frac{x}{\lambda}\right)^{n}.
\end{equation*}
事实上,由 $R(x, \lambda) - R(x, \lambda_{0}) = (\lambda - \lambda_{0})R(x, \lambda)R(x, \lambda_{0})$,以及当 $|\lambda| > \|x\|$ 时
\begin{equation*}
(x - \lambda e)^{-1} = \left(-\frac{1}{\lambda}\right)\left(e - \frac{x}{\lambda}\right)^{-1} = \left(-\frac{1}{\lambda}\right)\sum_{n=0}^{\infty}\left(\frac{x}{\lambda}\right)^{n}
\end{equation*}
即得.于是,考虑 $G(\lambda)$,若 $\Sp_{A}(x) = \emptyset$,则作为向量值函数 $G$ 在 $\mathbf{C}$ 上正则且有界.从而由向量值函数论(cf.\ \cite{ch3:3}),存在 $y \in A$ 使得 $G(\lambda) = y$($\lambda \in \mathbf{C}$).于是,$\lambda = 0$ 时 $y = G(0) = x^{-1}$,$\lambda = 1$ 时 $y = G(1) = (x - e)^{-1}$.这与 $y = G(\lambda) \in A^{-1}$ 相矛盾.因此 $\Sp_{A}(x) \neq \emptyset$.

当 $A$ 没有单位元时,由 $\Sp_{A}(x) = \Sp_{A_{e}}(x)$,把 $x$ 作为 $A_{e}$ 的元来考虑,由前半的论证可知 $\Sp_{A}(x)$ 是 $\mathbf{C}$ 的非空紧子集.
\end{proof}

\begin{corollary}\label{cor:3.1.7}
具有单位元 $e$ 的交换 Banach 代数 $A$ 若为域,则(作为环)与复数域同构.
\end{corollary}

\begin{proof}
设 $A$ 为域.若 $x \in A$,则由定理~\ref{thm:3.1.6},$\Sp_{A}(x) \neq \emptyset$.因此存在 $\lambda \in \mathbf{C}$ 使得 $x - \lambda e \notin A^{-1}$.此时因 $A$ 是域,$x - \lambda e = 0$.最终得到 $x = \lambda e$,由这一对应 $x \mapsto \lambda$,$A$ 与 $\mathbf{C}$ 作为环同构.
\end{proof}

\begin{definition}\label{def:3.1.8}
对 $x \in A$,定义
\begin{equation*}
\rho(x) = \sup\{|\lambda|:\ \lambda \in \Sp_{A}(x)\},
\end{equation*}
称 $\rho(x)$ 为 $x$ 的谱半径.
\end{definition}

\begin{theorem}\label{thm:3.1.9}
下式成立:
\begin{equation*}
\rho(x) = \lim_{n \to \infty}\|x^{n}\|^{1/n}.
\end{equation*}
\end{theorem}

\begin{proof}
设 $A$ 具有单位元 $e$.若 $\lambda \in \Sp_{A}(x)$,则 $\lambda^{n} \in \Sp_{A}(x^{n})$.由定理~\ref{thm:3.1.6},$|\lambda^{n}| \leq \|x^{n}\|$,故 $|\lambda| \leq \varliminf_{n \to \infty}\|x^{n}\|^{1/n}$.因此
\begin{equation*}
\rho(x) \leq \varliminf_{n \to \infty}\|x^{n}\|^{1/n}.
\end{equation*}
现在,考虑
\begin{equation*}
F(\lambda) = \left(-\frac{1}{\lambda}\right)\sum_{n=0}^{\infty}\left(\frac{x}{\lambda}\right)^{n},
\end{equation*}
当 $|\lambda| > \varlimsup_{n \to \infty}\sqrt[n]{\|x^{n}\|}$ 时,$F(\lambda)$ 的右边在 $A$ 中收敛,且 $F(\lambda) \in A$(cf.\ \cite{ch3:3}, \cite{ch3:8}).又,当 $\|\lambda^{-1}x\| < 1$ 时,由定理~\ref{thm:3.1.6},
\begin{equation*}
(x - \lambda e)^{-1} = F(\lambda),
\end{equation*}
故注意 $(x - \lambda e)F(1/\lambda) = e$.进一步,若 $|\lambda| > \rho(x)$,则由 $\rho(x)$ 的定义 $\lambda \notin \Sp_{A}(x)$,故存在正则的 $R(x, \lambda) = (x - \lambda e)^{-1}$.综上,
\begin{equation*}
\frac{1}{\rho(x)} \leq \frac{1}{\varlimsup_{n \to \infty}\sqrt[n]{\|x^{n}\|}}.
\end{equation*}
于是
\begin{equation*}
\varlimsup_{n \to \infty}\|x^{n}\|^{1/n} \leq \rho(x)
\end{equation*}
与前半合并即得
\begin{equation*}
\rho(x) = \lim_{n \to \infty}\|x^{n}\|^{1/n}.
\end{equation*}
当 $A$ 没有单位元时,把 $x \in A$ 放到 $A_{e}$ 中考虑,由 $\Sp_{A}(x) = \Sp_{A_{e}}(x)$ 归结为具有单位元的情形.
\end{proof}

\subsection{理想}\label{sec:3.1.3}

以下假设 $A$ 为交换 Banach 代数.

\begin{definition}\label{def:3.1.10}

(1) 若 $A$ 的子集 $I$ 是 $\mathbf{C}$ 上的环,且当 $x \in I$,$y \in A$ 时 $xy \in I$ 成立,则称 $I$ 为 $A$ 的理想.

(2) 若 $I$ 是理想,且不存在满足 $I \subsetneq J \subsetneq A$ 的理想 $J$,则称 $I$ 为极大理想.
\end{definition}

下面容易证明.

\begin{proposition}\label{prop:3.1.11}
当 $I$ 为 $A$ 的理想时,$I$ 的等价类 $A/I = \{x + I:\ x \in A\}$ 通过
(i) $(x + I) + (y + I) = x + y + I$,
(ii) $\lambda(x + I) = \lambda x + I$,
(iii) $(x + I)(y + I) = xy + I$
成为 $\mathbf{C}$ 上的环.$A/I$ 称为商环.特别地,当 $I$ 为闭集时,$A/I$ 由商拓扑 $\|x + I\|_{A/I} = \inf_{z \in I}\|x + z\|$ 成为交换 Banach 代数.
\end{proposition}

\begin{definition}\label{def:3.1.12}
$A$ 的理想 $I$ 称为正则理想,是指商环 $A/I$ 具有单位元.此时,$A/I$ 具有单位元的充要条件是存在 $u \in A$ 使得对任意 $x \in A$ 有 $ux - x \in I$.称这个 $u$ 为 $A$ 模 $I$ 的单位元.
\end{definition}

\begin{lemma}\label{lem:3.1.13}
当 $I \subsetneq A$ 且 $I$ 为正则理想时,存在正则极大理想 $J$ 满足 $I \subset J \subsetneq A$.
\end{lemma}

\begin{proof}
用 Zorn 引理证明.设 $I$ 为与 $A$ 不同的正则理想,则 $A$ 模 $I$ 的单位元 $u$ 显然不属于 $I$.事实上,若 $u \in I$,则由 $ux - x \in I$($x \in A$)得 $I = A$.现在,令 $\mathcal{J}$ 为满足 $I \subset J$,$u \notin J$ 的 $A$ 的理想 $J$ 的全体,则 $\mathcal{J}$ 关于包含关系构成半序集.取 $\mathcal{J}$ 的任意全序子集 $\mathcal{K}$,则 $L = \bigcup_{J \in \mathcal{K}} J$ 是包含 $I$ 的 $A$ 的理想,且 $u \notin L$.因此 $L \in \mathcal{J}$.从而由 Zorn 引理,$\mathcal{J}$ 中存在极大元 $J$.即满足 $I \subset J \subsetneq A$ 的理想.此时,对任意 $x \in A$,$ux - x \in I \subset J$,故 $J$ 是包含 $I$ 的正则极大理想.
\end{proof}

\begin{proposition}\label{prop:3.1.14}

(i) 若 $A$ 具有单位元 $e$,则 $A$ 的理想 $I$ 是正则理想,且 $A$ 的极大理想是闭集(是闭集的理想称为闭理想).

(ii) 设 $A$ 没有单位元.若 $I \subset A$ 是正则极大理想,则 $I$ 是闭理想.
\end{proposition}

\begin{proof}
(i) 当 $x + I \in A/I$ 时,由 $(e + I)(x + I) = x + I$,$I$ 是正则理想.设 $I$ 为极大理想.若 $\|e - x\| < 1$,则由定理~\ref{thm:3.1.3},$x \notin I$.因此若 $x \in I$,则 $\|e - x\| \geq 1$.于是用 $\overline{I}$ 表示 $I$ 的闭包,任取 $y \in \overline{I}$,则 $\|e - y\| \geq 1$.从而 $\overline{I} \subsetneqq A$,且 $\overline{I}$ 是理想.又因 $I$ 是极大理想,$I = \overline{I}$.因此是闭集.

(ii) 设 $u$ 为 $A$ 模 $I$ 的单位元.此时令 $\|u - I\| = \inf_{x \in I}\|u - x\|$,则 $\|u - I\| \geq 1$.事实上,假设存在 $v \in I$ 使得 $\|u - v\| < 1$.固定任意 $x \in A$,令
\begin{equation*}
y = x + \Bigl(\sum_{n=1}^{\infty}(u - v)^{n}\Bigr)x,
\end{equation*}
则
\begin{align*}
y - uy + vy & = x + \Bigl(\sum_{n=1}^{\infty}(u - v)^{n}\Bigr)x \\
& \quad - u\Bigl(x + \Bigl(\sum_{n=1}^{\infty}(u - v)^{n}\Bigr)x\Bigr) + v\Bigl(x + \Bigl(\sum_{n=1}^{\infty}(u - v)^{n}\Bigr)x\Bigr) \\
& = x + \Bigl(\sum_{n=1}^{\infty}(u - v)^{n}\Bigr)x - (u - v)x - \Bigl(\sum_{n=2}^{\infty}(u - v)^{n}\Bigr)x \\
& = x
\end{align*}
即,由 $v \in I$ 与 $u$ 的定义,$x \in I$.于是 $A = I$,矛盾.因此 $\|u - I\| \geq 1$ 成立.$\overline{I}$ 是 $I$ 的闭包,由 $I \subset \overline{I}$,$\overline{I}$ 是正则理想.因 $\|u - x\| \geq 1$($x \in I$),故 $\|u - x\| \geq 1$($x \in \overline{I}$).于是 $\overline{I} \subsetneqq A$.从而由 $I$ 是极大理想,$I = \overline{I}$.
\end{proof}

\begin{lemma}\label{lem:3.1.15}
当 $I$ 为与 $A$ 不同的正则理想时,$I$ 是正则极大理想的充要条件是 $A/I$ 与 $\mathbf{C}$ 同构.
\end{lemma}

\begin{proof}
设 $A/I$ 与 $\mathbf{C}$ 同构.若 $I$ 不是极大理想,则由补题~\ref{lem:3.1.13},存在极大理想 $J$ 使得 $I \subsetneq J$.取 $w \in J \cap I^{c}$,则与 $w + I$ 对应的是非零的 $\lambda \in \mathbf{C}$.因此 $w + I$ 在 $A/I$ 中具有逆元.于是 $J = A$,产生矛盾,故 $I$ 是极大理想.

反之,设 $I$ 为正则极大理想.由命题~\ref{prop:3.1.11},命题~\ref{prop:3.1.14},$A/I$ 成为具有单位元 $u + I$ 的交换 Banach 代数.此时 $A/I$ 是域.事实上,设 $x_{0} + I \in A/I$ 且 $x_{0} \notin I$,则由 $I$ 是极大理想,$\{x_{0}x + I:\ x \in A\} = A/I$.因为 $\{xx_{0} + y:\ x \in A,\ y \in I\}$ 是包含 $I$ 的理想.因此存在 $z \in A$ 使得 $zx_{0} + I = (z + I)(x_{0} + I) = u + I$,故 $x_{0} + I$ 具有逆元,$A/I$ 是域.从而由系~\ref{cor:3.1.7},$A/I$ 与 $\mathbf{C}$ 同构.
\end{proof}

\begin{theorem}\label{thm:3.1.16}
设 $A$ 为没有单位元的交换 Banach 代数,$A_{e}$ 为对 $A$ 添加单位元 $e$ 而得到的交换 Banach 代数.进一步,令 $M_{1}$ 为 $A_{e}$ 的正则极大理想的全体,$M_{0}$ 为 $A$ 的正则极大理想的全体,则 $M_{1} \setminus \{A\} \ni I \mapsto I \cap A \in M_{0}$ 是双射.
\end{theorem}

\begin{proof}
首先证明若 $I \in \mathfrak{M}_{1} \setminus \{A\}$,则 $I \cap A \in \mathfrak{M}_{0}$.首先,$I \cap A$ 是 $A$ 的理想,由假设容易看出 $I \not\subset A$.进一步,存在 $u \in A$ 使得 $u - e \in I$.事实上,因 $I \not\subset A$,存在非零复数 $\lambda_{0}$ 与 $y \in A$ 使得 $y + \lambda_{0}e \in I$.即 $y/\lambda_{0} + e \in I$.令 $-u = y/\lambda_{0}$,则 $e - u \in I$.由此得 $ux - x \in I$($x \in A$).因此 $ux - x \in I \cap A$($x \in A$),故 $I \cap A$ 是 $A$ 的正则理想.进一步,$I \cap A$ 是 $A$ 的极大理想.因为,若存在 $A$ 的理想 $J$ 满足 $A \cap I \subsetneq J \subset A$,则存在 $v_{0} \in J$ 使得 $v_{0} \notin I \cap A$.此时,$v_{0} + I \in A_{e}/I$ 由补题~\ref{lem:3.1.15} 对应于适当的非零复数 $\lambda' \in \mathbf{C}$,故 $v_{0}/\lambda' - e \in I$.于是由 $e - u \in I$,$v_{0}/\lambda' - e \in I$ 得 $v_{0}/\lambda' - u \in I \cap A \subset J$.再考虑 $v_{0} \in J$,得 $u \in J$.从而 $J = A$.综上得 $I \cap A \in \mathfrak{M}_{0}$.

其次证明上述映射是单射.事实上,设 $I_{1}, I_{2} \in \mathfrak{M}_{1} \setminus \{A\}$,$I_{1} \cap A = I_{2} \cap A$.此时,由前述,存在 $u \in A$ 使得 $ux - x \in I_{1} \cap A$($x \in A$).于是令
\begin{equation*}
J = \{x \in A_{e}:\ xu \in I_{1} \cap A\},
\end{equation*}
则 $J$ 是 $A_{e}$ 的理想且 $J \supset I_{1}$.这里,因 $I_{1} \cap A \subsetneq A$,$e \notin J$,得 $J \subsetneq A_{e}$.从而由 $I_{1}$ 是极大理想,$J = I_{1}$.同理 $J = I_{2}$.因此 $I_{1} = I_{2}$.

进一步证明该映射是满射.事实上,固定任意 $J \in M_{0}$,取适当的 $u \in A$ 使 $ux - x \in J$($x \in A$),令
\begin{equation*}
I' = \{x \in A_{e}:\ ux \in J\},
\end{equation*}
则 $I'$ 是 $A_{e}$ 的理想,且 $u \notin I'$.事实上,若 $u \in I'$,则 $u^{2} \in J$,由 $u^{2} - u \in J$ 得 $u \in J$,矛盾.于是取满足 $I \supset I'$ 且 $u \notin I$ 的 $A_{e}$ 的极大理想 $I$,由 $u \notin I$ 得 $I \neq A$,故 $I \in \mathfrak{M}_{1} \setminus \{A\}$.此时,若 $x \in J$,则 $ux \in J$,故 $x \in I'$.因此 $J \subset I$,即 $J \subset I \cap A$.因 $J$ 是极大理想,得 $J = I \cap A$.
\end{proof}

例．对任意 $x_{0} \in \mathbf{T}$,
\begin{equation*}
I(x_{0}) = \{f \in C(\mathbf{T}):\ f(x_{0}) = 0\}
\end{equation*}
是交换 Banach 代数 $C(\mathbf{T})$ 的极大理想.

\subsection{Gelfand 变换}\label{sec:3.1.4}

\begin{definition}\label{def:3.1.17}
取交换 Banach 代数 $A$,若 $\chi:\ A \to \mathbf{C}$ 具有如下性质,则称 $\chi$ 为特征.

(i) $\chi(\alpha x + \beta y) = \alpha\chi(x) + \beta\chi(y)$,\quad (ii) $\chi(xy) = \chi(x)\chi(y)$ \quad($x, y \in A$,$\alpha, \beta \in \mathbf{C}$)

用 $\Delta(A)$ 表示 $A$ 的满足 $\chi \neq 0$ 的特征 $\chi$ 的全体.
\end{definition}

\begin{theorem}\label{thm:3.1.18}
$\Delta(A)$ 与 $A$ 的正则极大理想全体 $\mathfrak{M}(A)$ 通过对应 $\Delta(A) \ni \chi \mapsto \chi^{-1}(0) \in \mathfrak{M}(A)$ 构成双射.
\end{theorem}

\begin{proof}
首先证明上述对应是满射.事实上,设 $I \in \mathfrak{M}(A)$,则由补题~\ref{lem:3.1.15},$A/I$ 与 $\mathbf{C}$ 同构,此时对应
\begin{equation*}
A/I \ni x + I \mapsto \chi(x + I) \in \mathbf{C}
\end{equation*}
是同构.进一步,令 $\chi(x) = \chi(x + I)$,则 $\chi$ 是特征.事实上,可定义性由若 $x - x' \in I$ 则按 $\chi$ 的定义 $\chi(x) = \chi(x + I) = \chi(x' + I)$ 可知.此时 $\chi^{-1}(0) = I$.反之,对 $\Delta(A) \ni \chi$,令 $I = \chi^{-1}(0)$,因 $\chi$ 是同态,$A/I$ 与 $\mathbf{C}$ 同构.

其次证明单射.事实上,若 $I = \chi^{-1}(0) = (\chi')^{-1}(0) = I'$,则 $\chi(x + I) = \chi'(x + I)$,故 $\chi = \chi'$.
\end{proof}

\begin{corollary}\label{cor:3.1.19}
设 $A$ 为交换 Banach 代数,则对 $x \in A$ 下式成立.

(1) 当 $A$ 具有单位元时,$\Sp_{A}(x) = \{\chi(x):\ \chi \in \Delta(A)\}$.

(2) 当 $A$ 没有单位元时,$\Sp_{A}(x) = \{\chi(x):\ \chi \in \Delta(A)\} \cup \{0\}$.
\end{corollary}

\begin{proof}
(1) 设 $A$ 的单位元为 $e$.当 $\lambda \in \Sp_{A}(x)$ 时,$x - \lambda e \notin A^{-1}$.此时存在 $I \in \mathfrak{M}(A)$ 使得 $x - \lambda e \in I$.从而由定理~\ref{thm:3.1.18},存在 $\chi \in \Delta(A)$ 使 $\chi^{-1}(0) = I$.即 $\chi(x - \lambda e) = 0$.于是 $\chi(x) = \lambda$.因此 $\lambda \in \{\chi(x):\ \chi \in \Delta(A)\}$.反向是显然的.

(2) 设 $A_{e}$ 为对 $A$ 添加单位元 $e$ 而得到的交换 Banach 代数,由 $\Sp_{A}(x) = \Sp_{A_{e}}(x)$,
\begin{equation*}
\Sp_{A_{e}}(x) = \{\chi(x):\ \chi \in \Delta(A_{e})\} = \{\chi(x):\ \chi \in \Delta(A)\} \cup \{0\}
\end{equation*}
即得.
\end{proof}

\begin{corollary}\label{cor:3.1.20}
当 $\chi \in \Delta(A)$ 时,$|\chi(x)| \leq \|x\|$($x \in A$).特别地,$\chi$ 是 $A$ 上的复值连续函数.
\end{corollary}

\begin{proof}
由系~\ref{cor:3.1.19},$\Sp_{A}(x) = \{\chi(x):\ \chi \in \Delta(A)\}$,或 $\Sp_{A}(x) = \{\chi(x):\ \chi \in \Delta(A)\} \cup \{0\}$.由于 $\rho(x) = \sup\{|\lambda|:\ \lambda \in \Sp_{A}(x)\}$,故 $|\chi(x)| \leq \|x\|$.特别地,当 $\|x_{n} - x\| \to 0$ 时,
\begin{equation*}
|\chi(x_{n}) - \chi(x)| = |\chi(x_{n} - x)| \leq \|x_{n} - x\| \to 0
\end{equation*}
故 $\chi:\ A \to \mathbf{C}$ 连续.
\end{proof}

\begin{definition}\label{def:3.1.21}
对 Banach 空间 $A$,把 $A$ 上的连续线性复值函数(称为 $A$ 上的连续线性泛函)的全体定义为 $A^{*}$.此时,$u \in A^{*}$ 在 $A$ 上连续的充要条件是存在常数 $C > 0$ 使得 $|u(x)| \leq C\|x\|$($x \in A$).$A^{*}$ 由通常的运算成为向量空间.
\end{definition}

\begin{definition}\label{def:3.1.22}

(1) 令 $u \in A^{*}$,$\|u\| = \sup\{|u(x)|:\ x \in A,\ \|x\| = 1\}$,则 $A^{*}$ 关于 $\|\cdot\|$ 成为 Banach 空间.称这个拓扑为 $A^{*}$ 的范数拓扑.

(2) 取 $u_{0} \in A^{*}$,把 $u_{0}$ 的邻域 $U$ 定义如下.
\begin{align*}
U &= U(u_{0};\ x_{1}, \ldots, x_{n}, \varepsilon) \\
&= \{u \in A^{*}:\ |u(x_{j}) - u_{0}(x_{j})| < \varepsilon\ (j = 1, \ldots, n)\}
\end{align*}
其中 $\varepsilon > 0$,$x_{1}, \ldots, x_{n} \in A$,$n \in \mathbf{N}$($\mathbf{N}$ 表示自然数集).此时,
\begin{equation*}
W(u_{0}) = \{U = U(u_{0};\ x_{1}, \ldots, x_{n}, \varepsilon):\ \varepsilon > 0,\ x_{1}, \ldots, x_{n} \in A,\ n \in \mathbf{N}\}
\end{equation*}
构成 $u_{0}$ 的基本邻域系.由它导出的 $A^{*}$ 的拓扑称为 $*$-弱拓扑.

(3) 在 $\Delta(A)$ 上引入 $A^{*}$ 的 $*$-弱拓扑的相对拓扑.
\end{definition}

\begin{lemma}\label{lem:3.1.23}

(1) $\Delta(A)$ 是 Hausdorff 空间.

(2) $\Delta(A)$ 在 $A^{*}$ 中关于 $*$-弱拓扑的闭包 $\overline{\Delta(A)}$ 为:当 $A$ 具有单位元时,$\overline{\Delta(A)} = \Delta(A)$;当 $A$ 没有单位元时,$\overline{\Delta(A)} = \Delta(A) \cup \{0\}$.其中 $0 \in A^{*}$ 为 $0(x) \equiv 0$($x \in A$).
\end{lemma}

\begin{proof}
(1) 设 $\chi_{1} \neq \chi_{2}$,$\chi_{1}, \chi_{2} \in \Delta(A)$.由 $\chi_{1} \neq \chi_{2}$,存在 $x \in A$ 使 $\chi_{1}(x) \neq \chi_{2}(x)$.于是令
\begin{equation*}
\varepsilon = \frac{|\chi_{1}(x) - \chi_{2}(x)|}{3} > 0.
\end{equation*}
进一步令 $U_{j} = \{\chi \in \Delta(A):\ |\chi(x) - \chi_{j}(x)| < \varepsilon\}$($j = 1, 2$),则 $U_{1} \cap U_{2} = \emptyset$.因为若 $U_{1} \cap U_{2} \ni \chi_{0}$,则 $|\chi_{0}(x) - \chi_{j}(x)| < \varepsilon$.于是 $|\chi_{1}(x) - \chi_{2}(x)| < 2\varepsilon = 2|\chi_{1}(x) - \chi_{2}(x)|/3$,矛盾.

(2) 取 $u_{0} \in \overline{\Delta(A)}$.对 $x, y \in A$,$\varepsilon > 0$,
\begin{equation*}
U = \{u \in A^{*}:\ |u(x) - u_{0}(x)| < \varepsilon,\ |u(y) - u_{0}(y)| < \varepsilon,\ |u(xy) - u_{0}(xy)| < \varepsilon\}
\end{equation*}
这是 $u_{0}$ 的邻域,故 $U \cap \Delta(A) \neq \emptyset$.因此存在 $\chi \in U \cap \Delta(A)$ 满足 $|\chi(x) - u_{0}(x)| < \varepsilon$,$|\chi(y) - u_{0}(y)| < \varepsilon$,$|\chi(xy) - u_{0}(xy)| < \varepsilon$.从而
\begin{align*}
& |u_{0}(xy) - u_{0}(x)u_{0}(y)| \leq |u_{0}(xy) - \chi(xy)| + |\chi(xy) - \chi(x)\chi(y)| \\
& \qquad + |\chi(y) - u_{0}(y)||\chi(x)| + |\chi(x) - u_{0}(x)||u_{0}(y)| \\
& \qquad \leq \varepsilon(1 + \|x\| + |u_{0}(y)|).
\end{align*}
因此证得 $u_{0}(xy) = u_{0}(x)u_{0}(y)$,故 $u_{0} \in \Delta(A) \cup \{0\}$.若 $A$ 具有单位元,则 $u_{0} \in \Delta(A)$,故 $\overline{\Delta(A)} = \Delta(A)$.没有单位元时,$\overline{\Delta(A)} = \Delta(A) \cup \{0\}$.
\end{proof}

下面是在函数分析中众所周知的定理.

\begin{theorem}\label{thm:3.1.24}
令 $B = \{u \in A^{*}:\ \|u\| \leq 1\}$,则 $B$ 在 $*$-弱拓扑下紧.
\end{theorem}

利用定理~\ref{thm:3.1.24} 可得下面结论.

\begin{corollary}\label{cor:3.1.25}
对交换 Banach 代数 $A$ 下式成立:

(1) 若 $A$ 具有单位元,则 $\Delta(A)$ 紧.

(2) 若 $A$ 没有单位元,则 $\Delta(A)$ 的每一点具有紧邻域.
\end{corollary}

\begin{proof}
(1) 由 $\overline{\Delta(A)} = \Delta(A) \subset B = \{u \in A^{*}:\ \|u\| \leq 1\}$,利用定理~\ref{thm:3.1.24} 可知 $\Delta(A)$ 紧.

(2) 由 $\overline{\Delta(A)} = \Delta(A) \cup \{0\}$($\subset B$),利用定理~\ref{thm:3.1.24} 可知 $\Delta(A) \cup \{0\}$ 紧.因此 $\Delta(A)$ 的每一点具有紧邻域.
\end{proof}

\begin{definition}\label{def:3.1.26}
设 $A$ 为交换 Banach 代数,$\mathfrak{M}(A)$ 为 $A$ 的正则极大理想的全体.按如下方式在 $\mathfrak{M}(A)$ 上引入拓扑:由定理~\ref{thm:3.1.18},$\mathfrak{M}(A)$ 与 $\Delta(A)$ 一一对应,故把 $\mathfrak{M}(A) \ni M$ 与 $\chi_{M} \in \Delta(A)$ 对应来考虑.取 $M_{0} \in \mathfrak{M}(A)$ 的邻域
\begin{equation*}
\begin{aligned}
U(M_{0}) &= \{M \in \mathfrak{M}(A):\ |\chi_{M}(x_{j}) - \chi_{M_{0}}(x_{j})| < \varepsilon\ (j = 1, \ldots, n)\} \\
&\bigl(= U(M_{0};\ \{x_{j}\}_{1}^{n}, \varepsilon)\bigr)
\end{aligned}
\end{equation*}
(其中 $x_{j} \in A$,$\varepsilon > 0$,$n \in \mathbf{N}$),令
\begin{equation*}
\mathcal{W}(M_{0}) = \{U(M_{0};\ \{x_{j}\}_{1}^{n}, \varepsilon):\ x_{j} \in A\ (j = 1, \ldots, n),\ \varepsilon > 0,\ n \in \mathbf{N}\}
\end{equation*}
为 $M_{0}$ 的基本邻域系.这样通过 $\Delta(A)$ 的拓扑而引入拓扑的 $\mathfrak{M}(A)$ 称为 $A$ 的极大理想空间.
\end{definition}

\begin{definition}\label{def:3.1.27}
设 $A$ 为交换 Banach 代数,$\mathfrak{M}(A)$ 为 $A$ 的极大理想空间.对 $x \in A$,令 $\hat{x}(M) = \chi_{M}(x)$($\chi_{M} \in \Delta(A)$),则 $\hat{x}$ 是 $\mathfrak{M}(A)$ 上的复值函数.称其为 $x$ 的 Gelfand 变换(Gelfand 表示).在不会引起混淆时,把 $\Delta(A)$ 与 $\mathfrak{M}(A)$ 等同,把 $\hat{x}(M)$ 表为 $\hat{x}(\chi_{M})$.
\end{definition}

\begin{proposition}\label{prop:3.1.28}

(1) $\hat{x}$ 是 $\mathfrak{M}(A)$ 上的连续函数.

(2) $\mathfrak{M}(A)$ 的拓扑是使 $\hat{x}$ 连续的最弱拓扑.

(3) 对任意 $\varepsilon > 0$,$\{M \in \mathfrak{M}(A):\ |\hat{x}(M)| \geq \varepsilon\}$ 是 $\mathfrak{M}(A)$ 的紧子集.
\end{proposition}

\begin{proof}
(1) 取 $x \in A$,$M_{0} \in \mathfrak{M}(A)$,$\varepsilon > 0$.令 $U = \{M \in \mathfrak{M}(A):\ |\hat{x}(M) - \hat{x}(M_{0})| < \varepsilon\}$,则由 $\mathfrak{M}(A)$ 的拓扑定义,$U$ 是 $M_{0} \in \mathfrak{M}(A)$ 的邻域.若 $M \in U$,则 $|\hat{x}(M) - \hat{x}(M_{0})| < \varepsilon$,故 $\hat{x}$ 在 $M_{0}$ 连续.因此 $\hat{x}$ 是 $\mathfrak{M}(A)$ 上的连续函数.

(2) 取 $x \in A$,$\varepsilon > 0$,$M_{0} \in \mathfrak{M}(A)$.令 $U = \{M \in \mathfrak{M}(A):\ |\hat{x}(M) - \hat{x}(M_{0})| < \varepsilon\}$.设 $\mathcal{O}$ 为集合 $\mathfrak{M}(A)$ 上使 $\hat{x}$ 连续的拓扑的开集系.此时,因 $\hat{x}$ 关于 $\mathcal{O}$ 的拓扑连续,$U \in \mathcal{O}$.因此 $A$ 的极大理想空间 $\mathfrak{M}(A)$ 的开集的全体包含于 $\mathcal{O}$.从而极大理想空间 $\mathfrak{M}(A)$ 的拓扑是使 $\hat{x}$ 连续的 $\mathfrak{M}(A)$ 的最弱拓扑.

(3) 由系~\ref{cor:3.1.25} 证明.
\end{proof}

\begin{definition}\label{def:3.1.29}
设 $A$ 为交换 Banach 代数.此时称 $A$ 为半单的,是指若 $x \in \bigcap_{M \in \mathfrak{M}(A)} M$,则 $x = 0$.
\end{definition}

下面的定理在函数分析中众所周知.

\textbf{闭图像定理}．X, Y 为 Banach 空间,$T:\ X \to Y$ 为线性算子.进一步,若当
\begin{equation*}
\lim_{n \to \infty} x_{n} = x, \quad \lim_{n \to \infty} Tx_{n} = y
\end{equation*}
时有 $Tx = y$,则 $T$ 连续.

\begin{theorem}\label{thm:3.1.30}
设 $A, B$ 为交换 Banach 代数,$B$ 半单.此时,若 $T:\ A \to B$ 是线性映射且 $T(xy) = (Tx)(Ty)$($x, y \in A$),则 $T$ 连续.
\end{theorem}

\begin{proof}
由闭图像定理,只需证明当 $x, x_{n} \in A$ 且 $z \in B$,
\begin{equation*}
\lim_{n \to \infty} x_{n} = x \quad \text{且} \quad \lim_{n \to \infty} Tx_{n} = z
\end{equation*}
时,$Tx = z$.对任意 $\chi \in \Delta(B)$,$\chi(T(xy)) = \chi(TxTy) = \chi(Tx)\chi(Ty) = (\chi \circ T)(x)(\chi \circ T)(y)$,故 $\chi \circ T$ 是 $A$ 的特征.因此 $\chi \circ T \in \Delta(A) \cup \{0\}$.从而 $\chi \circ T:\ A \to \mathbf{C}$ 连续.因此当 $\lim_{n \to \infty} x_{n} = x$ 时,
\begin{equation*}
\lim_{n \to \infty}\chi \circ T(x_{n}) = \chi \circ T(x).
\end{equation*}
另一方面,由 $\lim_{n \to \infty} Tx_{n} = z$,$\chi \in \Delta(B)$,
\begin{equation*}
\lim_{n \to \infty}\chi(Tx_{n}) = \chi(z),
\end{equation*}
故 $\chi \circ T(x) = \chi(z)$.即 $\chi_{M}(Tx) = \chi_{M}(z)$($M \in \mathfrak{M}(B)$).于是 $0 = \chi_{M}(Tx - z) = (Tx - z)^{\wedge}(M)$($\chi \in \Delta(B)$).因 $B$ 半单,$Tx = z$.
\end{proof}

\begin{corollary}\label{cor:3.1.31}
若 $A$ 是半单的交换 Banach 代数,则 $A$ 的范数拓扑唯一确定.
\end{corollary}

\begin{proof}
设 $(A, \|\cdot\|_{1})$,$(A, \|\cdot\|_{2})$ 为半单的交换 Banach 代数.在定理~\ref{thm:3.1.30} 中令 $T$ 为恒等映射.$T$ 线性且 $T(xy) = T(x)T(y)$,故 $T$ 满足定理~\ref{thm:3.1.30} 的假设.因此存在 $C_{1} > 0$ 使得 $\|x\|_{2} \leq C_{1}\|x\|_{1}$.反之同理,故存在 $C > 0$ 使得 $C^{-1}\|x\|_{1} \leq \|x\|_{2} \leq C\|x\|_{1}$($x \in A$).
\end{proof}

例．(1) 交换 Banach 代数 $C(\mathbf{T})$ 的 Gelfand 变换是恒等映射.事实上,用定理~\ref{thm:3.1.16} 后的例中的记号,极大理想具有 $I(x_{0})$($x_{0} \in \mathbf{T}$)的形式,且由 $x_{0} \in \mathbf{T}$ 到 $I(x_{0}) \in \mathfrak{M}(C(\mathbf{T}))$ 的映射连续,$\mathbf{T}$ 与 $\mathfrak{M}(C(\mathbf{T}))$ 同胚.因此成为同胚映射.

(2) 交换 Banach 代数 $\ell^{1}$ 的 Gelfand 变换对 $a = \{a_{n}\}$ 为 $\hat{a}(x) = \sum_{n \in \mathbf{Z}} a_{n}e^{inx}$.极大理想空间是 $\mathbf{T}$.事实上,由 $x \in \mathbf{T}$ 到 $M_{x} = \{a = \{a_{n}\} \in \ell^{1}:\ \hat{a}(x) = 0\} \in \mathfrak{M}(\ell^{1})$ 的映射是双射且连续.因此 $\mathbf{T}$ 与 $\mathfrak{M}(\ell^{1})$ 同胚(参照 3.2 节 $A(\mathbf{T})$ 一节).

\subsection{正则交换 Banach 代数}\label{sec:3.1.5}

\begin{definition}\label{def:3.1.32}
设 $A$ 为交换 Banach 代数.

(1) 对 $D \subset A$,令 $h(D) = \{M \in \mathfrak{M}(A):\ M \supset D\}$.此时利用 Gelfand 变换,$h(D) = \{\chi \in \Delta(A):\ \hat{x} = 0\ (x \in D)\}$.

(2) 对 $E \subset \mathfrak{M}(A)$,令 $k(E) = \bigcap\{M:\ M \in E\}$.此时利用 Gelfand 变换,也可认为 $k(E) = \{x \in A:\ \hat{x}(\chi) = 0\ (\chi \in E)\}$.其中,对空集 $\emptyset$,令 $k(\emptyset) = A$.
\end{definition}

\begin{theorem}\label{thm:3.1.33}
设 $E, E_{1}, E_{2} \subset \mathfrak{M}(A)$,则下式成立.

(1) $hk(E) \supset E$.

(2) $hk(E) = hk(hk(E))$.

(3) $hk(\emptyset) = \emptyset$.

(4) $hk(E_{1} \cup E_{2}) = hk(E_{1}) \cup hk(E_{2})$.
\end{theorem}

\begin{proof}
(1) 若 $\chi \in E$,则当 $x \in k(E)$ 时 $\hat{x}(\chi) = 0$,故 $\chi \in hk(E)$.

(2) 由 (1),$khk(E) \subset k(E)$.另一方面,对 $\chi \in hk(E)$ 取 $x \in k(E)$,由定义~\ref{def:3.1.32},$\hat{x}(\chi) = 0$,故 $k(E) \subset khk(E)$.因此 $khk(E) = k(E)$.从而得 $kh(hk(E)) = hk(E)$.

(3) 由 $k(\emptyset) = A$,$h(A) = \emptyset$,$hk(\emptyset) = \emptyset$.

(4) 由 $E_{1} \cup E_{2} \supset E_{j}$($j = 1, 2$),$hk(E_{1} \cup E_{2}) \supset hk(E_{j})$($j = 1, 2$).因此 $hk(E_{1} \cup E_{2}) \supset hk(E_{1}) \cup hk(E_{2})$.反之,设 $\chi \notin hk(E_{1}) \cup hk(E_{2})$,则存在 $x_{j} \in k(E_{j})$,$\chi(x_{j}) \neq 0$ 的 $x_{j} \in A$($j = 1, 2$).此时 $x_{1}x_{2} \in k(E_{1} \cup E_{2})$.事实上,若 $\chi' \in E_{1} \cup E_{2}$,当 $\chi' \in E_{1}$ 时,由 $x_{1} \in k(E_{1})$,$\chi'(x_{1}) = 0$,故 $\chi'(x_{1}x_{2}) = 0$.同理,$\chi' \in E_{2}$ 的情形也有 $\chi'(x_{1}x_{2}) = 0$.因此由 $\chi(x_{1}x_{2}) \neq 0$,$\chi \notin hk(E_{1} \cup E_{2})$.于是若 $\chi \notin hk(E_{1}) \cup hk(E_{2})$,则 $\chi \notin hk(E_{1} \cup E_{2})$,故 $hk(E_{1} \cup E_{2}) \subset hk(E_{1}) \cup hk(E_{2})$.
\end{proof}

\begin{definition}\label{def:3.1.34}
对 $E \subset \mathfrak{M}(A)$ 对应 $hk(E)$,由定理~\ref{thm:3.1.33} 在 $\mathfrak{M}(A)$ 上引入了闭包拓扑(把满足 $E = hk(E)$ 的 $E$ 作为闭集来引入拓扑).称其为 $\mathfrak{M}(A)$ 上的 hk 拓扑.
\end{definition}

\begin{definition}\label{def:3.1.35}
对交换 Banach 代数 $A$,令 $\hat{A} = \{\hat{x}:\ x \in A\}$.设 $F$ 为 $\mathfrak{M}(A)$ 的闭集,$\chi_{0} \notin F$,$\chi_{0} \in \mathfrak{M}(A)$,若存在 $x \in A$ 使得 $\hat{x}(\chi) = 0$($\chi \in F$)且 $\hat{x}(\chi_{0}) \neq 0$,则称交换 Banach 代数 $A$ 为正则的.
\end{definition}

\begin{theorem}\label{thm:3.1.36}
下列命题等价.

(1) 交换 Banach 代数 $A$ 正则.

(2) $\mathfrak{M}(A)$ 的拓扑与 hk 拓扑一致.

特别地,当 $A$ 具有单位元时,下列也等价.

(3) 设 $F$ 为 $\mathfrak{M}(A)$ 的闭集,$K$ 为 $\mathfrak{M}(A)$ 的紧集且 $K \cap F = \emptyset$,则存在 $x \in A$ 使得 $\hat{x}(\chi) = 0$($\chi \in F$),$\hat{x}(\chi) = 1$($\chi \in K$).
\end{theorem}

\begin{proof}
首先证明 (1) 与 (2) 等价.设 $F$ 为 $\mathfrak{M}(A)$ 的闭集.此时 $hk(F) \subset F$.事实上,由 (1) 的假设,若 $\chi \notin F$,则存在 $x \in A$ 使 $\hat{x}(F) = 0$,$\hat{x}(\chi) \neq 0$.此时由 $x \in k(F)$,$\chi \notin hk(F)$,故得 $hk(F) \subset F$.从而由定理~\ref{thm:3.1.33},$F \subset hk(F)$,故 $hk(F) = F$.因此当 $F$($\subset \mathfrak{M}(A)$)为闭集时,在 hk 拓扑下也是闭集.反之,设 $F$($\subset \mathfrak{M}(A)$)在 hk 拓扑下为闭集.此时 $hk(F) = F$,且 $h(k(F))$ 是 $\mathfrak{M}(A)$ 的闭集,故 $F$ 是 $\mathfrak{M}(A)$ 的闭集.假设 (3),则 (1) 显然.反之假设 (1) 证明 (3).为此考虑下列引理.
\end{proof}

\begin{lemma}\label{lem:3.1.37}
设 $A$ 为具有单位元 $e$ 的正则交换 Banach 代数,$I$ 为 $A$ 的理想,$E$ 为 $\mathfrak{M}(A)$ 的非空紧集且 $h(I) \cap E = \emptyset$.此时存在 $x \in I$ 使得 $\hat{x}(\chi) = 1$($\chi \in E$).
\end{lemma}

\begin{proof}
商环 $A/k(E)$ 的极大理想空间 $\mathfrak{M}(A/k(E))$ 与 $E$ 同胚.事实上,若 $\chi \in E$,令 $\widetilde{\chi}(x + k(E)) = \chi(x)$,则可定义 $A/k(E)$ 上的非零特征.进一步,把 $\chi \in E$ 对应到 $\widetilde{\chi}$ 的映射是从 $E$ 到 $\Delta(A/k(E))$ 的双射.因为,对 $\widetilde{\chi_{0}} \in \Delta(A/k(E))$,令 $\chi_{0}(x) = \widetilde{\chi_{0}}(x + k(E))$,若 $x \in k(E)$ 则 $\widetilde{\chi_{0}}(x) = 0$,由 (1),$hk(E) = E$,故 $\chi_{0} \in E$.因此该对应是满射.单射显然.又该映射连续,故 $E$ 与 $\mathfrak{M}(A/k(E))$ 同胚.现在,由假设 $h(I \cup k(E)) \subset h(I) \cap hk(E) = h(I) \cap E = \emptyset$,故 $h(I \cup k(E)) = \emptyset$.因此不存在同时包含 $I$,$k(E)$ 的 $A$ 的极大理想.即同时包含 $I$,$k(E)$ 的理想是 $A$.于是设 $e$ 为 $A$ 的单位元,存在 $u \in I$,$u \notin k(E)$ 使 $u + k(E) = e + k(E)$.因此存在 $v \in k(E)$ 使 $u = e + v$.即 $u \in I$,且 $\hat{u}(\chi) = \hat{e}(\chi) = 1$($\chi \in E$).
\end{proof}

((1) 蕴含 (3) 的证明)

在补题~\ref{lem:3.1.37} 中令 $I = k(F)$,$E = K$,则存在 $x \in A$ 使 $x \in k(F)$,$\hat{x}(\chi) = 1$($\chi \in K$).即 $x \in A$,$\hat{x}(\chi) = 0$($\chi \in F$),$\hat{x}(\chi) = 1$($\chi \in K$).

\begin{theorem}\label{thm:3.1.38}
设 $A$ 为具有单位元 $e$ 的正则交换 Banach 代数.设 $U_{1}, \ldots, U_{n}$ 为 $\mathfrak{M}(A)$ 的开覆盖.此时存在 $x_{j} \in A$($j = 1, \ldots, n$)使得 $\supp \widehat{x}_{j} \subset U_{j}$($j = 1, \ldots, n$),$\sum_{j=1}^{n} \widehat{x}_{j}(\chi) = 1$($\chi \in \mathfrak{M}(A)$).
\end{theorem}

\begin{proof}
$n = 1$ 时取 $x_{1} = e$ 即可.假设 $n - 1$ 时成立.设 $V, V_{3}, V_{4}, \ldots, V_{n}$ 为 $\mathfrak{M}(A)$ 的开覆盖且 $V \subset \overline{V} \subset U_{1} \cup U_{2}$,$V_{j} \subset \overline{V_{j}} \subset U_{j}$($j = 3, \ldots, n$).此时由 $n - 1$ 的假设,存在 $x, x_{3}, \ldots, x_{n} \in A$ 满足 $\supp \hat{x} \subset V$,$\supp \widehat{x_{j}} \subset V_{j}$,$\hat{x}(\chi) + \widehat{x_{3}}(\chi) + \cdots + \widehat{x_{n}}(\chi) = 1$($\chi \in \mathfrak{M}(A)$).此时令 $E = \supp \hat{x}$,$E_{1} = E \setminus U_{2}$,$E_{2} = E \setminus U_{1}$,由定理~\ref{thm:3.1.36} 存在 $y \in A$ 使 $\supp \hat{y} \subset V \cap U_{1}$,$\hat{y}(\chi) = 1$($\chi \in E_{1}$).事实上,$(V \cap U_{1})^{c} \cap E_{1} = \emptyset$.此时令 $x_{1} = xy$,$x_{2} = x(e - y)$,则 $x_{1} + x_{2} = x$,且 $\supp \widehat{x_{1}} \subset V \cap U_{1} \subset U_{1}$,$\supp \widehat{x_{2}} \subset E \cap U_{2} \subset U_{2}$.因此满足 $\widehat{x_{1}} + \cdots + \widehat{x_{n}} = 1$,$\supp \widehat{x_{j}} \subset U_{j}$($j = 1, \ldots, n$).
\end{proof}

\begin{corollary}\label{cor:3.1.39}
设 $A$ 为具有单位元 $e$ 的半单正则交换 Banach 代数.设 $I$ 为 $A$ 的理想,$f$ 为 $A$ 的极大理想空间 $\mathfrak{M}(A)$ 上的连续函数.此时,若对任意 $\chi_{0} \in \mathfrak{M}(A)$ 存在 $x_{0} \in I$ 以及 $\chi_{0}$ 的某邻域 $U_{0}$,使得 $f(\chi) = \widehat{x_{0}}(\chi)$($\chi \in U_{0}$),则存在 $x \in I$ 使 $\hat{x}(\chi) = f(\chi)$($\chi \in \mathfrak{M}(A)$).
\end{corollary}

\begin{proof}
对每个 $\chi_{0} \in \mathfrak{M}(A)$ 按假设对应 $\chi_{0}$ 的邻域 $U_{0}$,则 $\{U_{0}\}$ 是 $\mathfrak{M}(A)$ 的开覆盖,因 $\mathfrak{M}(A)$ 紧,从 $\{U_{0}\}$ 中可取 $\mathfrak{M}(A)$ 的有限开覆盖 $U_{1}, \ldots, U_{n}$.对此,由定理~\ref{thm:3.1.38} 选取 $z_{1}, \ldots, z_{n} \in A$,$\supp \widehat{z}_{j} \subset U_{j}$,$\widehat{z}_{1} + \cdots + \widehat{z}_{n} = 1$.此时
\begin{equation*}
f(\chi)\widehat{z_{1}}(\chi) + \cdots + f(\chi)\widehat{z_{n}}(\chi) = f(\chi) \quad (\chi \in \mathfrak{M}(A))
\end{equation*}
成立.这里,由假设,不妨设存在 $x_{j} \in I$ 使 $f(\chi) = \hat{x}_{j}(\chi)$($\chi \in U_{j}$),故
\begin{equation*}
\widehat{x_{1}}(\chi)\widehat{z_{1}}(\chi) + \cdots + \widehat{x_{n}}(\chi)\widehat{z_{n}}(\chi) = f(\chi) \quad (\chi \in \mathfrak{M}(A))
\end{equation*}
成立.因 $A$ 半单且 $I$ 是理想,令 $x = x_{1}z_{1} + \cdots + x_{n}z_{n}$,则 $x \in I$,$\hat{x} = f$.
\end{proof}

例．(1) 交换 Banach 代数 $C(\mathbf{T})$ 的 Gelfand 变换是恒等映射,故正则.

(2) 交换 Banach 代数 $\ell^{1}$ 的极大理想空间是 $\mathbf{T}$,且 $\mathbf{T}$ 上的 2 次连续可微函数 $f(x)$ 可表为 $f(x) = \sum_{n \in \mathbf{Z}} a_{n}e^{inx}$,$\sum_{n \in \mathbf{Z}} |a_{n}| < \infty$,故 $\ell^{1}$ 正则(参照 3.2.1 小节).

\subsection{谱综合}\label{sec:3.1.6}

\begin{definition}\label{def:3.1.40}
设 $A$ 为正则交换 Banach 代数.

(1) 设 $E \subset \mathfrak{M}(A)$ 为闭集,$I$ 为正则闭理想.当满足 $h(I) = E$ 的 $I$ 唯一确定时,称 $E$ 为交换 Banach 代数 $A$ 的谱综合集.

(2) 当 $\mathfrak{M}(A)$ 中存在不是谱综合集的闭集时,称交换 Banach 代数 $A$ 为谱综合不可能的.

(3) 当 $\mathfrak{M}(A)$ 的所有闭集都是谱综合集时,称交换 Banach 代数 $A$ 为谱综合可能的.
\end{definition}

以下在本节中,设 $A$ 为具有单位元 $e$ 的正则交换 Banach 代数.

\begin{definition}\label{def:3.1.41}
对 $E \subset \mathfrak{M}(A)$,令 $I(E) = \{x \in A:\ \hat{x}(\chi) = 0\ (\chi \in E)\}$,则 $I(E)$ 是闭理想.又令 $j(E) = \{x \in A:\ \supp \hat{x} \text{ 是紧集,且 } \supp \hat{x} \cap E = \emptyset\}$,$J(E) = \overline{j(E)}$($j(E)$ 在 $A$ 中的闭包),则 $j(E), J(E)$ 都是理想.
\end{definition}

\begin{lemma}\label{lem:3.1.42}
设 $A$ 半单.若 $I \subset \mathfrak{M}(A)$ 为闭集,则 $I(E)$ 是满足 $h(I) = E$ 的理想 $I$ 中最大的理想,$j(E)$ 是满足 $h(I) = E$ 的理想 $I$ 中最小的理想.
\end{lemma}

\begin{proof}
当 $h(I) = E$ 时,若 $x \in I$,则 $\widehat{x}(\chi) = 0$($\chi \in E$),故 $x \in I(E)$.又 $h(j(E)) \supset E$ 显然.若 $\chi \notin E$,则因 $A$ 正则,存在 $x_{0} \in A$ 使得 $\widehat{x_{0}}(\chi) = 1$,$\supp \widehat{x_{0}}$ 是紧集,$\supp \widehat{x_{0}} \cap E = \emptyset$.此时 $x_{0} \in j(E)$,故 $\chi \notin h(j(E))$,得 $h(j(E)) = E$.又取满足 $h(I) = E$ 的理想 $I$.此时,对满足 $K \cap E = \emptyset$ 的紧集 $K \subset \mathfrak{M}(A)$,由补题~\ref{lem:3.1.37} 存在 $a \in I$ 使 $\widehat{a}(\chi) = 1$($\chi \in K$),$\widehat{a}(\chi) = 0$($\chi \in E$).因此,对满足 $\supp \widehat{x} \subset K$ 的 $x \in A$,$\widehat{a}\widehat{x} = \widehat{x}$,因 $A$ 半单,$x = ax \in I$.从而 $j(E) \subset I$.
\end{proof}

\begin{note}\label{note:3.1.43}
由补题~\ref{lem:3.1.42},闭集 $E \subset \mathfrak{M}(A)$ 是 $A$ 的谱综合集的充要条件是 $I(E) = J(E)$.
\end{note}

\begin{definition}\label{def:3.1.44}
交换 Banach 代数 $A$ 满足 Ditkin 条件,是指当 $\chi \in \mathfrak{M}(A)$ 与 $x \in A$ 满足 $\hat{x}(\chi) = 0$ 时,存在 $\{x_{n}\} \subset A$ 使得 $\widehat{x_{n}} = 0$ 在 $\chi$ 的某邻域内成立,且 $\lim_{n \to \infty} x_{n}x = x$.
\end{definition}

\begin{theorem}\label{thm:3.1.45}
设 $A$ 为具有单位元 $e$ 的半单正则交换 Banach 代数,且满足 Ditkin 条件.设 $I$ 为 $A$ 的闭理想,$x_{0} \in kh(I)$.此时,若 $h(\{x_{0}\}) \cap h(I)$ 的拓扑边界不含非空完全集,则 $x_{0} \in I$.
\end{theorem}

\begin{proof}
对给定的 $x_{0} \in kh(I)$,令
\begin{equation*}
E = \{\chi \in \Delta(A):\ \text{对 } I \text{ 的任意元 } x \text{,不存在 } \chi \text{ 的邻域使 } \widehat{x_{0}} = \widehat{x} \text{ 成立}\},
\end{equation*}
若能证明 $E = \emptyset$,则由系~\ref{cor:3.1.39} 即得定理~\ref{thm:3.1.45}.现在,若 $\chi \notin h(I)$,由补题~\ref{lem:3.1.37},存在 $a \in I$ 及 $\chi$ 的邻域使 $\widehat{a} = 1$,故 $\chi \notin E$.又,若 $\chi \in \mathfrak{M}(A)$ 不在 $h(\{x_{0}\})$ 的拓扑边界上,则 $\chi$ 是 $h(\{x_{0}\})$ 的内点或 $\chi \notin h(\{x_{0}\})$.若 $\chi$ 是 $h(\{x_{0}\})$ 的内点,则 $\widehat{x_{0}} = 0$ 在 $\chi$ 的某邻域内成立,故 $\chi \notin E$.当 $\chi \notin h(\{x_{0}\})$ 时,由 $x_{0}$($\in kh(I)$)的假设,$h(\{x_{0}\}) \supset hkh(I) = h(I)$,故 $\chi \notin h(I)$.因此 $\chi \notin E$.综上,证明了 $E$ 包含于 $h(\{x_{0}\}) \cap h(I)$ 的拓扑边界.进一步,$E$ 没有孤立点.事实上,若 $\chi \in E$ 是孤立点,则存在 $\chi$ 的邻域 $U$,使得 $U$ 在 $\mathfrak{M}(A)$ 中的闭包 $\overline{U}$ 满足 $\overline{U} \cap E = \{\chi\}$.因 $\widehat{x_{0}}(\chi) = 0$,由 Ditkin 条件,存在 $\{z_{n}\} \subset A$ 使得 $\widehat{z_{n}} = 0$ 在 $\chi$ 的邻域内成立,且 $\lim_{n \to \infty} z_{n}x_{0} = x_{0}$.此时,取在 $\chi$ 的邻域内 $\widehat{y} = 1$ 且 $\supp \widehat{y} \subset U$ 的 $y \in A$,由系~\ref{cor:3.1.39},$z_{n}yx_{0} \in I$.因此由 $\lim_{n \to \infty} z_{n}yx_{0} = yx_{0} \in I$,$yx_{0} \in I$.于是由 $E$ 的定义,$\chi \notin E$,产生矛盾.综上,$E$ 是完全集.由假设,$h(\{x_{0}\}) \cap h(I)$ 的拓扑边界不含非空完全集,故 $E = \emptyset$.
\end{proof}

\begin{corollary}\label{cor:3.1.46}
设 $A$ 为具有单位元 $e$ 的半单正则且满足 Ditkin 条件的交换 Banach 代数,若闭集 $E \subset \mathfrak{M}(A)$ 的拓扑边界不含非空完全集,则 $E$ 是($A$ 的)谱综合集.
\end{corollary}

\begin{proof}
设 $x_{0} \in I(E)$.此时,由 $h(J(E)) = E$,$x_{0} \in k(E) = kh(J(E))$.进一步,利用假设,$h(\{x_{0}\}) \cap E = E$,故 $h(\{x_{0}\}) \cap h(J(E))$ 的拓扑边界不含完全集.因此由定理~\ref{thm:3.1.45},$x_{0} \in J(E)$,故证得 $I(E) = J(E)$.
\end{proof}

\subsection{半单交换 Banach 代数的作用函数}\label{sec:3.1.7}

\begin{definition}\label{def:3.1.47}
设 $A$ 为半单交换 Banach 代数.设 $G \subset \mathbf{C}$,$F:\ G \to \mathbf{C}$ 为函数.对 $x \in A$,$\Sp_{A}(x) \subset G$,若存在 $y \in A$ 使得
\begin{equation*}
F(\widehat{x}(M)) = \widehat{y}(M) \quad (M \in \mathfrak{M}(A))
\end{equation*}
则称 $F$ 为 $A$ 的作用函数.
\end{definition}

\begin{theorem}\label{thm:3.1.48}
设 $A$ 为具有单位元 $e$ 的半单交换 Banach 代数.若 $F$ 是单连通区域 $G$($\subset \mathbf{C}$)上的正则函数,则 $F$ 是 $A$ 的作用函数.即,设 $\Gamma$ 为 $G$ 中具有长度的闭曲线,且 $\supp \hat{x}$ 包含于 $\Gamma$ 的内部,则
\begin{equation*}
F(x) = \frac{1}{2\pi i}\int_{\Gamma}\frac{F(z)}{z - x}\dd z \in A
\end{equation*}
可以定义,且 $\widehat{F(x)}(M) = F(\widehat{x}(M))$($M \in \mathfrak{M}(A)$)成立.
\end{theorem}

\begin{proof}
由向量值函数的函数论证明(cf.\ \cite{ch3:3}, \cite{ch3:8}).细节从略.
\end{proof}

关于解析函数作用于 Banach 代数,在 \cite{ch3:1}, \cite{ch3:2} 等中有论述.
